% Compile with:
%   pdflatex paper && bibtex paper && pdflatex paper && pdflatex paper
\documentclass[11pt]{article}

\usepackage[a4paper, margin=1in]{geometry}
\usepackage{amsmath,amssymb}
\usepackage{graphicx}
\usepackage{booktabs}
\usepackage{caption}
\usepackage{subcaption}
\usepackage{natbib}
\usepackage{xcolor}
\usepackage{listings}
\usepackage[hidelinks]{hyperref}
\usepackage[final,protrusion=true,expansion=false]{microtype}

\graphicspath{%
  {../../out/figures/phase/}%
  {../../out/figures/scenarios/}%
  {../../out/figures/sweeps/}%
  {../../out/figures/policies/}%
  {../../out/figures/wealth_tax/}%
  {../../out/figures/scale_up/}%
}

\lstset{
  basicstyle=\ttfamily\footnotesize,
  frame=single,
  breaklines=true,
  columns=fullflexible,
  keepspaces=true,
}

\title{Unequal Agents:\\
\large An autocatalytic-closure agent-based model of wealth inequality,\\
cultural transmission, and the design of progressive wealth taxes}

\author{[Author]\\
\small (with supervision and acknowledgements TBD)}
\date{Working paper, draft of \today}

\begin{document}
\maketitle

\begin{abstract}
We present an agent-based model of wealth inequality in which seven
occupationally heterogeneous agent classes (workers, makers, employers,
lenders, rentiers, speculators, and a state) interact through production,
extraction, credit, taxation, inheritance, and a bandwidth-limited
cultural-transmission channel coupled to parent wealth.

The model supports two methodological contributions.  We introduce a
per-tick scalar \emph{productive-flow share}
$\Phi_t = F^\text{prod}_t / (F^\text{prod}_t + F^\text{ext}_t)$, defined
symmetrically across productive (wages, state transfers, maker
value-add, employer markup) and extractive (rent, interest, capital
return) accrual, and we explicitly distinguish
$\Phi$ from the topological autocatalytic-set closure of
\citet{kauffman1988evolution} and \citet{hordijk2004detecting} which it
takes its conceptual lineage from. We then couple cultural-transmission
\emph{bandwidth} to wealth on a log scale, so the wealth coupling stays
graded across orders of magnitude rather than saturating at the
upper-middle threshold as in a prior version of the model.

Across thirteen scenarios, five parameter sweeps, a clean
(tax $\times$ spending $\times$ inheritance) factorial, and a focused
suite of eleven real-world wealth-tax proposals (Zucman 2024, Piketty,
Warren, Sanders, Norway, Spain, UK 2010--2024 baseline, a Reform-flavour
deregulatory prospectus), we report eight findings.

(i) Three regimes (\emph{viable autocatalytic}, \emph{extractive
equilibrium}, \emph{collapse}) emerge robustly under k-means clustering
with silhouette 0.55. (ii) Spending and tax design and inheritance
taxation play distinct roles when factorialised: spending raises alive
counts by 40\% without moving $\Phi$; tax schedule compresses
household-sector inequality; inheritance taxation is the single most
$\Phi$-effective lever. (iii) Cultural transmission alone does not
equalise wealth even when fully equalised at birth, but a class-wise
skill-to-wealth elasticity regression shows the asymmetry is partly
mechanism: workers and makers have $R^2 \approx 0.9$ skill-to-wealth
fits, rentiers and lenders have $R^2 \leq 0.35$.  (iv) Inheritance
reform produces population collapse unless paired with a safety net.
(v) Zucman's flat 2\% above \pounds10M is not separable from a
progressive 2/4/8\% form at 10 seeds at $N=10{,}000$, and an analytical
fixed-point derivation $w^* = \tau \theta / (\tau - r)$ explains why:
the lower bracket sets the post-tax cap before higher tiers can fire.
(vi) The original draft's ``regressive trap'' headline is retracted:
under state-excluded top-share metrics top-1\% share decreases
monotonically with the wealth-tax rate at every threshold; the
inverted-monotonicity finding was entirely a state-contamination
artefact, and the genuine policy claim is the opposite: a wealth tax
at the \pounds10M threshold compresses top-1\% share by a factor of
4.4 in our $N=10{,}000$ runs.
(vii) At $N=10{,}000$ the UK Tory baseline produces 19 billionaires; the
Reform-flavour prospectus produces 97; every wealth-tax policy collapses
the count to 1, robustly across 10 seeds. (viii) A partial calibration
to the UK Sunday Times Rich List Pareto slope of $-0.706$ is achievable
with $r=0.06$, rent share 0.50, inheritance tax 0.40 (slope $-0.673$);
the levels of the model's tail remain too flat at the very top and too
fat at the lower-billionaire band by about an order of magnitude in
either direction, identifying a calibration target for follow-up work.

We also include a minimal open-ended-evolution demonstration of skill
discovery to address the related concern that the model's closed K=5
skill space is incompatible with the spirit of the alife literature.

This version of the paper substantially revises four claims from the
earlier draft following referee criticism, which the lead author asked
this revision to confront head-on: the State agent is now excluded from
all percentile and Gini statistics; $\Phi$ is defined symmetrically
across productive and extractive accrual; the regressive-trap claim is
fully retracted under state-exclusion; and the top-tier-does-not-bind
result is now derived analytically rather than reported as an empirical
equivalence.
\end{abstract}

%==========================================================================
\section{Introduction}
%==========================================================================

The empirical case that wealth concentration drives, rather than reflects,
economic dysfunction has sharpened over the past decade. The wealth-tax
proposal advanced by \citet{zucman2024global} for the G20 frames this
explicitly: existing tax instruments do not capture the very rich, and
the existing distribution is therefore self-reinforcing. Stevenson's
public articulation of the same claim, that inequality is a driver,
not a symptom, gives a political register but the same underlying
intuition.

We approach this question with a stylised agent-based simulation. The
model is not predictive. It is a stock-flow accounting machine populated
by occupationally heterogeneous agents whose interactions can be steered
by a panel of policy ``dials''. The point is to make visible, in motion,
the qualitative consequences of fiscal designs, and to identify where
they fail in unintuitive ways.

We organise the contribution around two ideas. The first is a
quantitative one. We define an \emph{autocatalytic closure} ratio
$\Phi_t \in [0,1]$ which is the share of total monetary flow at time $t$
that occurs between productive agent classes (workers, makers,
employers), versus between productive and extractive classes (rentiers,
lenders, speculators). The metric operationalises in a continuous,
per-tick form the productive-vs-extractive distinction that runs through
\citet{kauffman1988evolution}, \citet{mazzucato2018value}, and, less
formally, the entire heterodox tradition.
\citet{hordijk2023emergence} use a topological detection of autocatalytic
structure in a production-function graph; our $\Phi_t$ is a flow-based
sibling, and we believe it has not previously been operationalised in
this scalar form for an ongoing simulated economy.

The second contribution is mechanical. We give every agent a $K=5$
dimensional skill vector. When an agent dies and is replaced by a
child of the same occupational type, the child inherits a number of
skill dimensions that scales with the parent's terminal wealth. Cultural
transmission is therefore \emph{bandwidth-limited} in a way that depends
on economic position, a wealth-coupled fidelity mechanism that, to
our reading of the literature \citep{cavallisforza1981cultural,
henrich2004demography, lewis2012transmission, henrich2001evolution}, has
not been articulated in this form in any published agent-based model.

The remainder of the paper proceeds as follows. Section~\ref{sec:related}
positions the work. Section~\ref{sec:model} specifies the agents,
mechanics, metrics, and the wealth-tax instrument. Section~\ref{sec:method}
describes scenarios, parameter sweeps, and the eleven real-world policies
modelled. Section~\ref{sec:results} reports findings in order of
political-economic salience, including a scale-up to $N=10{,}000$ agents
where the billionaire class first becomes statistically representable.
Section~\ref{sec:discussion} discusses limitations, venue positioning,
and what we cannot model.

%==========================================================================
\section{Related work}
\label{sec:related}
%==========================================================================

The closest single paper to ours is \citet{spangler2022distribution},
which couples estate taxation, skill inheritance, and a Carnegie
effect within an ABM, finding that skill inheritance raises inequality
while the Carnegie effect lowers everyone's wealth. Our model differs
in three respects: (i) we add explicit extractive channels (rent,
capital return, interest) alongside the inheritance mechanism;
(ii) we couple skill-transmission \emph{bandwidth} (not target choice
or binary skill carry-over) to wealth; and (iii) we measure system
viability via the autocatalytic-closure ratio rather than terminal-period
Gini alone.

\citet{cardoso2025wealth}, recently published, makes a thematically
close point in a yard-sale-style kinetic ABM: \emph{social protection}
parameters dominate growth-driven redistribution. We recover an analogous
finding for state spending versus state revenue (\S\ref{sec:state-spending}),
and our wealth-tax sweeps (\S\ref{sec:threshold-rate}) further
disaggregate the social-protection result into a threshold-vs-rate
factorisation that, to our knowledge, is novel.

Within the econophysics tradition, the closest theoretical neighbours
are \citet{bouchaud2000wealth} on multiplicative wealth condensation,
\citet{yakovenko2009colloquium} on the two-class distribution, and
\citet{boghosian2017oligarchy} on the yard-sale model. These provide
the formal backdrop against which our extraction-channel mechanism
operates; we add a state, a credit sector, and a wealth-coupled
transmission layer.

\citet{piketty2014capital} provides the empirical and theoretical case
for $r > g$ dynamics, which the model recovers as an emergent property
under realistic-looking parameter settings (Figure~\ref{fig:phase}).
\citet{piketty2013theory} on optimal inheritance taxation is the
analytical version of our inheritance-trapdoor finding
(\S\ref{sec:inheritance}). \citet{berman2016modeling} present a related
ABM of US wealth concentration since 1980; \citet{pagliarini2017agent}
shows the JEIC venue's openness to calibrated wealth-distribution ABMs.

In cultural evolution, the bandwidth idea sits adjacent to
\citet{boyd1985culture}, \citet{henrich2015secret}, and especially
\citet{lewis2012transmission}'s ``fidelity is the key'' result. Where
\citet{henrich2004demography} ties cumulative culture to population
size, and \citet{henrich2001evolution} couples transmission target to
prestige, we couple transmission \emph{bandwidth} to wealth. We have
not found published precedent for that specific coupling.

%==========================================================================
\section{The model}
\label{sec:model}
%==========================================================================

\subsection{Agent types and population}

We model seven agent classes with fixed proportions per 100 agents:
72~workers, 9~makers, 7~employers, 5~lenders, 3~rentiers, 3~speculators,
and 1~state. The proportions are calibrated by hand to current
advanced-economy occupational structure, not from any one data source.

Each agent carries: wealth $w$, debt $d$, age $a$, lifespan
$\ell$ drawn at birth from $\mathcal{N}(70,12)$ truncated to $[20,120]$,
and a length-5 skill vector
$\mathbf{s} \in \mathbb{R}^5$ with dimensions interpreted as
\emph{technical}, \emph{organisational}, \emph{social}, \emph{financial},
\emph{generative}. The class of an agent biases its initial skill
distribution: makers receive a boost on generative and technical;
employers on organisational and financial; rentiers/lenders/speculators
on financial only; workers on social and technical.

For the scale-up experiments (\S\ref{sec:scale-up}) we introduce a
\emph{population\_scale} multiplier that scales agent counts and initial
wealth totals linearly while keeping the state a singleton and absolute
spending dials proportional to scale. Per-capita wealth statistics are
invariant under this scaling.

\subsection{Per-tick mechanics}

A single tick (one stylised year) runs the following phases in order.

\paragraph{Production.} Workers are sorted by skill and assigned to
makers and employers in greedy order until hiring capacities are
exhausted. Makers transform a material input plus paid labour into
output of value
$(\text{inputs}) \cdot M (0.7 + 0.6 s_5^{\text{maker}})$
where $M$ is the maker multiplier dial and $s_5$ is the maker's
generative skill. Employers pay wages and capture a markup
$(1 + \mu)$ on their workers' produced value. The
state optionally hires a configurable fraction of remaining unemployed
workers and purchases a budgeted volume of maker output, both
parameterised through dials.

\paragraph{Extraction.} Each worker pays a fixed rent
$\rho \cdot W$ which is distributed across rentiers in proportion to
rentier wealth. Each rentier earns a capital return $r$ on its existing
wealth. Each speculator's wealth evolves geometrically with drift
$\mu_s$ and volatility $\sigma_s$.

\paragraph{Consumption and credit.} Every non-state agent pays
subsistence $\sigma$ per tick. Agents with insufficient wealth borrow
from the aggregate lender pool, with the borrowed amount allocated to
lenders proportionally to their available capacity. If aggregate
capacity is exhausted the agent dies (``starves'').

\paragraph{Interest, repayment, write-off.} Indebted agents pay interest
at rate $i$ when liquid; otherwise interest is capitalised. Agents
beyond a configurable debt ceiling are written off; the loss is
distributed across lenders proportionally to wealth.

\paragraph{Lender deposits and recapitalisation.} A fixed fraction of
surplus wealth held by rentiers, employers, and makers is deposited
into the lender pool, modelling the macro-flow of financial savings.
The state optionally backstops the lender pool to a configurable floor.

\paragraph{Tax and UBI.} Every non-state living agent with positive
wealth pays
\begin{equation}
T(w) \;=\; \tau_0 w \;+\; \tau_1 \max(0, w - \theta_0)
       \;+\; \sum_{k=1}^{K} \tau^\text{tier}_k \max(0, w - \theta^\text{tier}_k)
\end{equation}
where $\tau_0$ is the flat rate, $\tau_1$ is the single-tier
progressive rate above threshold $\theta_0$, and the tiered marginal
rates $\tau^\text{tier}_k$ apply above their thresholds in a
\emph{stacked} form: a rate above a high threshold adds to (rather
than replacing) the rates above the lower thresholds. The state
collects $T$ and disburses universal basic income $U$ per agent to
the extent the state's accumulated wealth allows.

\paragraph{Reproduction and transmission.} Agents whose age exceeds
their lifespan die a natural death. A child of the same class is
created. Wealth passes after an inheritance tax $\iota$ and is capped
at $\bar{w}_{\text{inh}}$; revenue above the cap accrues to the state.
The child's skill vector is initialised by selecting the
\textsc{bandwidth} most-skilled dimensions of the parent's vector,
where the bandwidth in dimensions is
\begin{equation}
B \;=\; \min\!\Big( K,\; \max\!\big(1,\; \lfloor b_0 K + b_w \cdot
(w_{\text{parent}}/100) \rfloor\big)\Big)
\end{equation}
with $b_0$ a base proportion and $b_w$ a wealth-coupling slope.
\emph{Knowledge loss} is recorded as the sum of skill mass present in
the parent but not transmitted.

\subsection{The autocatalytic closure metric}
\label{sec:phi}

During each tick we accumulate two flows treated symmetrically as
\emph{productive-class accrual} and \emph{extractive-class accrual}.
Productive accrual $F^\text{prod}_t$ sums (i) wages paid by makers,
employers, and the state to workers; (ii) state purchases of maker
output; (iii) maker value-add, i.e. the surplus $\text{output} -
\text{inputs}$ created in production and retained by the maker; and
(iv) the share of worker output captured by employers as markup that
contributes to the metabolic rate. Of the gross employer markup
($\text{worker\_output} \cdot \text{employer\_markup}$), a fraction
$0.4$ is treated as productive value-add (the organisational
contribution of the employer that creates the markup) and the
remaining $0.6$ as zero-sum capture from worker output; this
split is a modelling choice and we report sensitivity to it in
the dial documentation (Appendix~\ref{app:dials}).
Extractive accrual $F^\text{ext}_t$
sums rent paid by workers to rentiers, interest paid by debtors to
lenders, capital returns accruing to rentier stocks, and positive
speculator returns (the negative speculator returns are not
counted, on the same convention by which rent and interest do not
have negative counterparts). Both sums
mix inter-class transfers with intra-class value creation in the same
manner: a maker creating $0.6$ units of value per tick contributes to
$F^\text{prod}$ exactly as a rentier compounding at $r=0.04$ on its
existing stock contributes to $F^\text{ext}$. The autocatalytic
closure is
\begin{equation}
\Phi_t \;=\; \frac{F^\text{prod}_t}{F^\text{prod}_t + F^\text{ext}_t} \in [0,1].
\end{equation}
A value near 1 indicates that monetary accrual at time $t$ is
predominantly going to or being created by productive-class agents; a
value near 0 indicates that accrual is dominated by extractive
channels.

A point of nomenclature. $\Phi_t$ is
\emph{not} an autocatalytic-set closure in the technical sense of
\citet{kauffman1988evolution} or \citet{hordijk2004detecting}, who
detect a topological RAF property of a chemical or production graph.
$\Phi_t$ is a continuous, scalar, per-tick \emph{flow-share} measure
which operationalises in measurable form the
productive-vs-extractive distinction that runs through
\citet{kauffman1988evolution}, \citet{mazzucato2018value}, and the
classical political-economic tradition. We have retained the
``autocatalytic'' adjective because the original conceptual lineage is
the one we engage with; readers wanting a strictly graph-theoretic
quantity should consult Hordijk and colleagues instead. An alternative
name we considered, ``productive-flow share,'' is more honest about
what is measured and we use the two interchangeably where it aids
reading.

\subsection{Wealth-tax instrument}

For the policy experiments in \S\ref{sec:policy} we add the tiered
wealth-tax mechanism described above and adopt the convention that
\emph{one model wealth unit corresponds to \pounds100{,}000} of real
wealth. Reference thresholds in model units are: 100 ($=\pounds$10M,
Zucman's lower bracket), 500 ($=\pounds$50M, Warren), 1000
($=\pounds$100M), 10000 ($=\pounds$1B, the billionaire threshold).
This mapping is a convenience for labelling and not a calibration;
no parameter has been fitted to data.

\subsection{Scale}

We run all 100-agent experiments at \emph{population\_scale} = 1 and
the scale-up experiments at \emph{population\_scale} = 100, yielding
$N=10{,}000$. The implementation refactors four hot loops from
$O(N^2)$ to $O(N)$ by accumulating per-class totals and distributing
in single passes, and prunes dead agents from the agent list every
ten ticks. On seven processor cores a 200-tick run at $N=10{,}000$
completes in approximately 25 seconds; a full eight-policy comparison
with three seeds and a $7\times 7$ wealth-tax sweep with two seeds
completes in approximately ten minutes.

%==========================================================================
\section{Method}
\label{sec:method}
%==========================================================================

\subsection{Scenarios}

We define thirteen named scenarios at $N=100$, each a particular setting
of the dial panel. \texttt{current\_world} is the baseline; the others
span: an extreme-extraction regime (\texttt{stevenson\_world}); a
production-heavy regime (\texttt{maker\_economy}); two state-spending
regimes (\texttt{tax\_and\_spend}, \texttt{nordic}); two memory-mechanism
contrasts (\texttt{memory\_equalised}, \texttt{memory\_starved}); two
inheritance-reform regimes (\texttt{inheritance\_break},
\texttt{inheritance\_break\_ubi}); and four pathological or
policy-limit regimes (\texttt{feudal}, \texttt{creditless},
\texttt{pure\_speculation}, \texttt{socialist}). Each scenario is run
for 250 ticks with eight random seeds; we report mean trajectories
and inter-quartile ribbons.

\subsection{Parameter sweeps}

We perform five two-dimensional sweeps at $N=100$ on $9\times9$ grids
with three seeds each: \texttt{rent\_vs\_r}, \texttt{tax\_vs\_ubi},
\texttt{memory}, \texttt{state\_action}, and
\texttt{inheritance\_x\_ubi}. Each sweep varies two dials over the
ranges listed in Appendix~\ref{app:sweeps} while holding all others at
defaults. Three additional wealth-tax-specific sweeps (Sweep~A: single
tier with threshold $\times$ rate; Sweep~B: stacked-marginal lower-rate
$\times$ top-rate; the scale-up sweep at $N=10{,}000$) are introduced
in \S\ref{sec:threshold-rate} and \S\ref{sec:scale-up}.

\subsection{Policy scenarios}
\label{sec:policy-defs}

For the wealth-tax analysis we encode eleven real-world or
contemporaneous-political wealth-tax proposals. Each policy combines a
tax schedule with the spending package its political project advances,
because as our state-spending finding makes clear (\S\ref{sec:state-spending}),
the two cannot be evaluated in isolation. Policy specifications are
listed in Appendix~\ref{app:policies}. Effective-rate curves across
policies are shown in Figure~\ref{fig:erc}.

\begin{figure}[t]
\centering
\includegraphics[width=0.9\linewidth]{effective_rate_curves.png}
\caption{Effective annual wealth-tax rate as a function of wealth, by
policy. Vertical lines mark the \pounds10M, \pounds100M, and \pounds1B
thresholds. Most policies' progressive structures are visible as
piecewise-linear segments on the log-wealth axis. Zucman's flat
2\% above \pounds10M and his progressive 2/4/8\% version differ only
above \pounds100M, where the marginal additions stack.}
\label{fig:erc}
\end{figure}

%==========================================================================
\section{Results}
\label{sec:results}
%==========================================================================

We report results in increasing political-economic salience, ending
with the scale-up to $N=10{,}000$ where the billionaire class first
becomes representable.

\subsection{Regimes: a phase portrait}
\label{sec:phase}

Across the thirteen scenarios at $N=100$, scenarios cluster into three
clearly separated regions of the three-dimensional state space
(Gini, $\Phi_t$, alive). Figure~\ref{fig:phase} shows the trajectories
and Table~\ref{tab:regimes} summarises endpoint coordinates.

\begin{figure}[t]
\centering
\begin{subfigure}{0.55\linewidth}
\includegraphics[width=\linewidth]{phase3d_trajectories.png}
\end{subfigure}\hfill
\begin{subfigure}{0.42\linewidth}
\includegraphics[width=\linewidth]{phase3d_shadows.png}
\end{subfigure}
\caption{Phase portrait of all scenarios in the three-dimensional state
space (Gini, autocatalytic closure $\Phi$, alive). Three regions are
visible: viable autocatalytic ($\Phi > 0.4$, alive $> 65$); extractive
equilibrium ($\Phi \approx 0$, alive 40--75); collapse ($\Phi = 0$,
alive $< 20$).}
\label{fig:phase}
\end{figure}

\begin{table}[t]
\centering
\small
\caption{Endpoint coordinates by scenario at $N=100$, 250 ticks, mean
across 8 seeds, with the State agent \emph{excluded} from Gini and
share statistics (revision per cross-cutting review note). Regime
assignment is by k-means with $k=3$ on z-normalised
$(\text{Gini}, \Phi, \text{alive})$, silhouette score 0.55. Cluster
labels follow the centroid coordinates: viable (centroid $\Phi=0.75$,
alive 85), extractive equilibrium ($\Phi=0.01$, alive 62), collapse
($\Phi=0.00$, alive 12).}
\label{tab:regimes}
\begin{tabular}{lrrrl}
\toprule
scenario & Gini & $\Phi$ & alive & regime \\
\midrule
socialist             & 0.30 & 0.880 & 99.0 & viable autocatalytic \\
nordic                & 0.30 & 0.701 & 89.9 & viable autocatalytic \\
inheritance\_break\_ubi & 0.57 & 0.667 & 85.9 & viable autocatalytic \\
maker\_economy        & 0.77 & 0.586 & 68.8 & viable autocatalytic \\
tax\_and\_spend       & 0.43 & 0.567 & 83.6 & viable autocatalytic \\
memory\_equalised     & 0.92 & 0.015 & 76.9 & extractive equilibrium \\
creditless            & 0.87 & 0.013 & 55.2 & extractive equilibrium \\
current\_world        & 0.89 & 0.011 & 63.1 & extractive equilibrium \\
memory\_starved       & 0.87 & 0.011 & 55.5 & extractive equilibrium \\
pure\_speculation     & 0.92 & 0.001 & 57.9 & extractive equilibrium \\
stevenson\_world      & 0.88 & 0.001 & 20.2 & collapse \\
feudal                & 0.78 & 0.000 & 15.2 & collapse \\
inheritance\_break    & 0.00 & 0.000 &  0.0 & collapse \\
\bottomrule
\end{tabular}
\end{table}

\paragraph{State exclusion changes the picture.} The two largest
Gini and top-share changes between this revision and the original
draft are for the spending-heavy scenarios: \texttt{nordic}'s Gini
falls from 0.79 (with State) to 0.29 (without), \texttt{socialist}'s
from 0.81 to 0.30. The State was the single richest agent in those
runs; it accumulated tax revenue without dispersing it back to the
household sector and so dominated the percentile statistics. The
corrected numbers make the qualitative claim sharper: viable
state-spending regimes do not merely contain Gini, they actively
reduce household-sector inequality below 0.3, comparable to lived
Scandinavian distributions.

The $k=3$ clustering replaces an earlier ``by inspection'' regime
labelling. Silhouette scores are 0.55 for $k=2$, 0.55 for $k=3$,
0.42 for $k=4$, and 0.32 for $k=5$. The data does not statistically
prefer $k=3$ over $k=2$; the choice of three regimes is a
theoretical/interpretive decision motivated by the distinction
between extractive equilibrium and collapse, which we view as
meaningfully different policy outcomes even where the data only
weakly separates them. With 13 scenario points, any clustering
claim is fragile. The regimes are interpretable as three rather
than dictated as three, with the collapse-vs-extractive boundary
the least sharp. The viable regime is autocatalytic in the
productive-flow share sense defined in \S\ref{sec:phi} and is
reached by any of three policy mixes: heavy state spending, high
inheritance taxation, or elevating maker productivity.

\subsection{What state spending actually does: a clean factorial}
\label{sec:state-spending}

State purchases of maker output enter $F^\text{prod}$
directly (see \S\ref{sec:phi}), so toggling them on raises $\Phi$
by construction. To separate the definitional contribution from
the policy-relevant one, we run a clean factorial over three policy
levers (tax schedule, spending package, inheritance taxation),
reporting how each one moves the four headline metrics separately.

The factorial varies the tax schedule (none, Zucman flat 2\% above
threshold $\theta_1$, Zucman progressive 2/4/8\%), the spending
package (off, on at UBI 0.4 + state employment 0.3 + state purchases
at scaled rate 5), and the inheritance-tax rate ($\iota \in \{0,
0.4, 0.8\}$). Eighteen cells, six seeds, 200 ticks.
Table~\ref{tab:factorial} reports marginal means across each
factor.

\begin{table}[t]
\centering
\small
\caption{Marginal effects from the (tax $\times$ spending $\times$
inheritance) factorial, $N=100$, 6 seeds, 200 ticks. Each row
averages the median of the four marginalised cells.}
\label{tab:factorial}
\begin{tabular}{llrrrr}
\toprule
factor & level & alive & Gini & $\Phi$ & top-10 \\
\midrule
tax schedule & none                   & 66.5 & 0.85 & 0.186 & 0.77 \\
             & Zucman flat 2\%        & 67.8 & 0.68 & 0.449 & 0.56 \\
             & Zucman prog 2/4/8\%    & 67.8 & 0.67 & 0.465 & 0.55 \\
\midrule
spending     & off                    & 56.1 & 0.79 & 0.369 & 0.63 \\
             & on                     & 78.6 & 0.71 & 0.364 & 0.62 \\
\midrule
inheritance  & $\iota = 0.0$          & 68.8 & 0.74 & 0.234 & 0.65 \\
             & $\iota = 0.4$          & 67.5 & 0.74 & 0.337 & 0.66 \\
             & $\iota = 0.8$          & 65.8 & 0.74 & 0.529 & 0.56 \\
\bottomrule
\end{tabular}
\end{table}

Three findings drop out cleanly.

\paragraph{Spending saves lives but barely moves $\Phi$.} Turning the
spending package on raises alive from 56 to 79 (a 40\% increase) but
moves $\Phi$ by less than $0.005$ in either direction (0.369 vs
0.364). Under the symmetric $\Phi$ definition, the wage flow added by
state spending is offset, in share terms, by the reduced rentier
extraction; spending's headline effect is on mortality, not on the
flow composition. The original draft conflated these by counting
state transfers in the numerator without crediting maker value-add,
which artificially inflated the spending effect on $\Phi$. The
honest claim, restated, is that spending is a life-saving instrument
that operates orthogonally to $\Phi$.

\paragraph{Tax schedule reduces inequality and raises $\Phi$.} The
Zucman 2\% bracket above threshold $\theta_1$ moves Gini from 0.85 to
0.68 and top-10 share from 0.77 to 0.56. It raises $\Phi$ from 0.186
to 0.449, because the tax shrinks rentier stocks and therefore the
capital-return component of $F^\text{ext}$. The Zucman progressive
schedule adds nothing detectable on top of the flat 2\% in this
factorial, consistent with \S\ref{sec:top-tier-no-bind}.

\paragraph{Inheritance tax is the most $\Phi$-effective single
lever.} Going from $\iota = 0$ to $\iota = 0.8$ raises $\Phi$ from
0.23 to 0.53 (a $+0.30$ swing) while moving Gini only modestly. The
mechanism is the same as for the wealth tax: estate taxation
truncates dynastic wealth accumulation, shrinking the rentier stocks
that drive capital-return extraction. \citet{piketty2013theory}'s
analytical case for inheritance taxation is recovered as the
single most $\Phi$-effective single-lever intervention in our model.

\begin{figure}[t]
\centering
\includegraphics[width=0.95\linewidth]{tax_vs_ubi.png}
\caption{The \texttt{tax\_vs\_ubi} sweep. Eight outcome metrics on the
(flat\_tax\_rate, UBI) plane, with a state-spending baseline.
Re-reading post-revision: the alive surface (top left) shows the
genuine life-saving effect of UBI; $\Phi$ (top right) is dominated by
the tax-progressivity axis rather than by UBI, consistent with
Table~\ref{tab:factorial}.}
\label{fig:taxubi}
\end{figure}

So the revised summary of \S\ref{sec:state-spending} is that
\emph{three} levers each do something different and they should not
be sold as a single policy stance: spending keeps people alive, tax
schedules compress household inequality, and inheritance taxes raise
the productive-flow share. The original draft's ``spending flips the
regime'' headline conflated these and overclaimed what spending does
on its own.

\subsection{Visualising the flow regime: Sankey of steady-state flows}
\label{sec:sankey}

Figure~\ref{fig:sankey} shows mean per-tick wealth flows in steady
state for three policies, rendered as ribbon Sankey diagrams. Each
ribbon's width is proportional to the per-tick flow magnitude from a
source class to a destination class; node-box height is proportional to
total outflow (left) or inflow (right) of that class.

\begin{figure}[t]
\centering
\includegraphics[width=\linewidth]{sankey_flows.png}
\caption{Steady-state wealth flows for three policies at $N=10{,}000$,
mean across 3 seeds, ticks 80--180. \textbf{Left (uk\_tory\_2010\_24):}
44\% of total flow is the RENTIER--LENDER deposit loop; productive
wage flow (MAKER+EMPLOYER to WORKER) is 11\%. The low $\Phi = 0.05$
is visible as the small share of edges landing on WORKER.
\textbf{Middle (zucman\_2pct\_spend):} the
RENTIER--LENDER loop has shrunk to 14\%; wealth-tax revenue
(MAKER, RENTIER, LENDER to STATE) is 40\% of total flow; STATE to
WORKER appears at 7\%. \textbf{Right (socialist):}
WORKER--STATE--WORKER is the dominant circulation, accounting for
40\% of total flow; the RENTIER--LENDER loop has effectively
disappeared. Socialist $\Phi$ measured here is 0.88
(Table~\ref{tab:regimes}).}
\label{fig:sankey}
\end{figure}

The figure makes the productive-flow-share metric visible as a
diagram rather than a scalar. The Tory panel's low $\Phi$ is the
visual fact that most ribbons connect non-productive classes; the
Socialist panel's high $\Phi$ is the fact that almost every ribbon
ends at WORKER or MAKER. This Sankey is the cleanest single image
for explaining the project's central metric to a non-technical
audience.

A secondary observation worth noting. Total flow per tick under
\texttt{zucman\_2pct\_spend} (32{,}760 units) is slightly lower than
under \texttt{uk\_tory\_2010\_24} (35{,}788 units) and
\texttt{socialist} (35{,}203 units). The Zucman policy reduces total
monetary throughput while shifting its composition strongly toward
productive flows, a result consistent with the cardoso et al.\
\citeyearpar{cardoso2025wealth} finding that social-protection
parameters dominate growth-driven redistribution in kinetic-tax ABMs.

\subsection{Inheritance reform as a conditional trapdoor}
\label{sec:inheritance}

Figure~\ref{fig:inh} shows the \texttt{inheritance\_x\_ubi} sweep. At
the corner $(\iota=1, U=0)$, 100\% inheritance tax with no safety
net, only 2 of 100 agents survive 200 ticks. Children are born with
zero wealth and starve. At $(\iota=0.875, U=1.5)$, nearly equal
inheritance reform combined with strong UBI, 95 of 100 survive.
The two dials are not substitutes; inheritance reform is conditional
on the existence of an alternative wealth floor.

\begin{figure}[t]
\centering
\includegraphics[width=0.95\linewidth]{inheritance_x_ubi.png}
\caption{The \texttt{inheritance\_x\_ubi} sweep. The alive surface
(top-left) shows the trapdoor at $(\iota=1, U=0)$. The Gini surface
(top-right) shows the mortality-blindness paradox: lowest Gini
coincides with lowest alive.}
\label{fig:inh}
\end{figure}

The same plot also shows our \emph{mortality-blind Gini} paradox.
Lowest measured Gini sits in the corner where mortality is highest:
extreme extraction kills the bottom of the distribution, leaving a
homogeneous surviving population that registers as ``equal''. In
\texttt{rent\_vs\_r} (Figure~\ref{fig:rentr}) the same pattern
appears: highest Gini is at moderate rent and high $r$, while the
extreme-rent corner shows a falsely lower Gini. We compute
mortality-corrected Gini and top-share metrics that credit recent
deaths' peak lifetime wealth back into the denominator
(\S\ref{sec:phi}, \S\ref{sec:regressive-trap}).

Figure~\ref{fig:survival} shows the temporal shape of the
mortality dynamics behind these endpoint numbers.
\texttt{inheritance\_break} loses its population within the first
30 ticks; \texttt{feudal} and \texttt{stevenson\_world} follow
shortly after. \texttt{nordic} and \texttt{socialist} hold near
full population indefinitely; \texttt{inheritance\_break\_ubi}
maintains the bulk of the population once the UBI starts catching
the orphans.

\begin{figure}[t]
\centering
\includegraphics[width=0.95\linewidth]{survival_scenarios.png}
\caption{Survival curves for eight illustrative scenarios at
$N=100$, 6 seeds. \textbf{Left:} alive over time (count).
\textbf{Right:} cumulative excess mortality (starvation + debt
write-off) as fraction of initial population. Collapse scenarios
finish within 100 ticks of the start; viable scenarios hold a
near-stationary alive count throughout.}
\label{fig:survival}
\end{figure}

\begin{figure}[t]
\centering
\includegraphics[width=0.95\linewidth]{rent_vs_r.png}
\caption{The \texttt{rent\_vs\_r} sweep. Alive (top-left) drops
sharply with extraction; closure (bottom-left) is highest in the
gentle corner; Gini (top-right) is non-monotone in extraction because
of the mortality bias.}
\label{fig:rentr}
\end{figure}

\subsection{Cultural erosion under economic pressure}
\label{sec:culture}

The model exhibits cultural \emph{erosion} rather than cumulative
cultural evolution. The K=5 skill space is closed by construction:
mutation can reset a dimension to a fresh draw but no genuinely new
skill dimension can appear. The \citet{henrich2015secret}
``cumulative culture'' framing therefore does not strictly apply;
what the model demonstrates is the \emph{loss} of cultural breadth
under wealth concentration, in the spirit of the demographic model
of \citet{henrich2004demography}. \citet{vaesen2016cumulative}
critique the Tasmanian inference in that paper by arguing that the
relevant population sizes were not as different as Henrich's model
required, weakening the population-size-explains-loss conclusion;
our model assumes the loss mechanism (wealth-coupled bandwidth)
rather than testing it. A genuinely open-ended treatment of cultural
novelty, in the adjacent-possible spirit of
\citet{hordijk2023emergence}, is demonstrated separately in
\S\ref{sec:open-skills}.

\paragraph{Bandwidth coupling.} The bandwidth coupling
\begin{equation}
\text{bw\_dims} \;=\; b_0 K + b_w K \log_{10}\!\big(\max(1, w_{\text{parent}})\big) / 4
\end{equation}
clipped to $[1, K]$ stays graded across orders of magnitude of
wealth rather than saturating at the upper-middle threshold (as a
prior linear formulation did).
Figure~\ref{fig:bandwidth} decomposes the \texttt{memory} sweep
breadth response into the marginal effect of $b_0$ and of $b_w$. The
$b_0$ axis dominates (range 2.99) with the wealth-coupling
$b_w$ contributing a graded response (range 0.87) across the
full wealth distribution. The wealth coupling is a real, secondary
channel rather than a saturating step function.

\begin{figure}[t]
\centering
\includegraphics[width=0.95\linewidth]{bandwidth_decomp.png}
\caption{Decomposition of the \texttt{memory} sweep. \textbf{Left:}
cultural breadth on the $(b_0, b_w)$ plane. \textbf{Middle:}
marginal effect of $b_0$ (avg over $b_w$) and of $b_w$ (avg over
$b_0$); base bandwidth dominates, wealth coupling adds a graded
correction. \textbf{Right:} realised distribution of
$\text{bw\_dims}$ at birth by class under \texttt{current\_world}.
Higher-wealth classes (rentier, employer) routinely inherit 4--5
dimensions; workers cluster at 2--3.}
\label{fig:bandwidth}
\end{figure}

\paragraph{The asymmetry result, with the skill-wealth channel
quantified.} The original draft's negative result, that equalising
cultural transmission at birth does not equalise wealth, recurs
under the redesigned coupling. The \texttt{memory} sweep moves
cultural breadth between 1.3 and 5.0 (a fourfold range) while Gini
(now with State excluded) moves between 0.84 and 0.92, a 10\%
range. \texttt{memory\_equalised} gives the highest non-extractive
Gini (0.93); cultural equality alone does not move wealth concentration.

To distinguish structural from emergent asymmetry we quantify the
skill$\to$wealth channel directly.
Table~\ref{tab:elasticity} reports per-class OLS coefficients and
normalised elasticities for the regression of 40-tick wealth gain
on the skill vector, with $n=809$ agent-windows across 8 seeds at
$N=500$.

\begin{table}[t]
\centering
\small
\caption{Skill$\to$wealth normalised elasticities by class. Wealth
gain over 40 ticks regressed on the five-dim skill vector at
$N=500$. Normalised elasticity =
$\text{coef}\cdot\sigma_{\text{skill}}/|\overline{\text{gain}}|$.}
\label{tab:elasticity}
\begin{tabular}{lrrrrrrl}
\toprule
class & tech & org & social & fin & gen & $R^2$ & n \\
\midrule
WORKER     & 0.76 & 0.95 & $-0.02$ & $-0.03$ & $-0.03$ & 0.97 & 395 \\
MAKER      & 0.00 & 0.00 & 0.01 & 0.00 & 0.33 & 0.98 & 144 \\
EMPLOYER   & 0.30 & 0.03 & 0.25 & $-0.10$ & 0.24 & 0.12 & 100 \\
LENDER     & $-0.01$ & 0.00 & $-0.01$ & 0.00 & $-0.01$ & 0.13 & 79 \\
RENTIER    & $-0.06$ & $-0.05$ & 0.09 & $-0.10$ & $-0.08$ & 0.12 & 47 \\
SPECULATOR & 0.56 & $-0.39$ & $-1.41$ & 0.04 & 0.13 & 0.11 & 44 \\
\bottomrule
\end{tabular}
\end{table}

The channel is class-specific. Workers and makers have strong
class-typed skill$\to$wealth elasticities (worker:
technical~0.76, organisational~0.95, $R^2$~0.97; maker:
generative~0.33, $R^2$~0.98). Employers have moderate
elasticities ($R^2$~0.12) concentrated on technical and social.
Lenders, rentiers, and speculators all have $R^2 \in [0.11, 0.13]$
and small normalised elasticities ($|\eta| \leq 0.10$ for lenders
and rentiers). The asymmetry is therefore partly structural: in
the rentier and lender layers, capital-return and interest dynamics
dominate and skills contribute almost nothing. But in the productive
layers, skills do real work, and improving them raises per-tick
wealth gain. The cultural-economic asymmetry visible in the
\texttt{memory} sweep is the consequence of one half of the economy
(extractive classes) having weak internal skill-wealth feedback while
the other half (productive classes) has strong feedback.

\begin{figure}[t]
\centering
\includegraphics[width=0.95\linewidth]{memory.png}
\caption{The \texttt{memory} sweep. Cultural breadth (top middle)
varies by $4\times$; Gini (top left) varies by 10\%. Cultural
concentration responds to the bandwidth coupling but wealth
concentration is anchored by the extractive-channel dynamics.}
\label{fig:mem}
\end{figure}

\paragraph{Comparison to Spangler \& Sarkar.} The closest
ABM precedent, \citet{spangler2022distribution}, find that skill
inheritance raises wealth inequality. We can test that mechanism
directly in our model. A clean toggle holds the bandwidth-gating
mechanism and class-bias prior intact, and swaps only whether the
transmitted skill values come from the parent (the default) or from
a fresh class-typical Beta draw (the comparison arm); a coarser
toggle, which destroys the bandwidth mechanism along with parental
transmission by setting the mutation rate $\mu = 1.0$, is reported
for contrast. At $N=100$, 10 seeds, 250 ticks:

\begin{table}[h]
\centering
\small
\begin{tabular}{lrrr}
\toprule
metric & default & class-prior & no-inheritance \\
       & (parent-derived) & (clean toggle) & ($\mu=1$, confounded) \\
\midrule
Gini median   & 0.885 & 0.876 & 0.908 \\
Top-1 median  & 0.235 & 0.250 & 0.254 \\
Top-10 median & 0.785 & 0.683 & 0.891 \\
Alive median  & 61.5  & 56.5  & 73.0 \\
Breadth       & 3.38  & 3.55  & 5.00 \\
\bottomrule
\end{tabular}
\end{table}

Under the clean toggle, switching transmitted values from
parent-derived to class-typical-prior \emph{lowers} Gini by 0.010
and \emph{raises} top-1 share by 0.016. So parent-derived skill
inheritance does raise Gini relative to a class-typical prior. The
direction is a partial reproduction of Spangler \& Sarkar; the
magnitude is small. The coarser $\mu = 1$ arm gives the opposite
Gini direction (0.908 vs 0.885) because it simultaneously removes
the bandwidth gate and the class-bias persistence that the clean
toggle preserves.

The elasticity result (Table~\ref{tab:elasticity}) explains the
small magnitude. Workers and makers have strong skill-to-wealth
fits and their inherited skills do feed wealth; rentiers and
lenders have weak fits and their inherited skills are not the
binding constraint on their wealth dynamics. The Spangler effect
should therefore be smaller in a multi-channel ABM than in their
single-channel design.

\subsection{Wealth-tax policies in steady state}
\label{sec:policy}

Table~\ref{tab:policies-100} reports steady-state outcomes for the
eleven policy scenarios at $N=100$, 6 seeds, 250 ticks; trajectories
are shown in Figure~\ref{fig:policies-compare}.

\begin{table}[t]
\centering
\small
\caption{Eleven real-world wealth-tax proposals in steady state.
$N=100$, 6 seeds, 250 ticks. The Reform-flavour row is included as a
near-future political scenario rather than a present policy.}
\label{tab:policies-100}
\begin{tabular}{lrrrrrr}
\toprule
policy & alive & $\Phi$ & top10 & state $w$ & Gini & metabolic \\
\midrule
uk\_tory\_2010\_24            & 57 & 0.050 & 0.85 & 36{,}541 & 0.90 & 39.0 \\
uk\_labour\_2024              & 57 & 0.230 & 0.75 & 28{,}455 & 0.85 & 39.5 \\
zucman\_2pct\_hoard           & 51 & 0.418 & 0.59 & 17{,}943 & 0.75 & 39.1 \\
zucman\_2pct\_spend           & 80 & 0.412 & 0.54 & 25{,}215 & 0.58 & 43.1 \\
zucman\_progressive\_2\_4\_8 & 80 & 0.420 & 0.54 & 24{,}949 & 0.58 & 43.1 \\
warren\_2020                  & 77 & 0.264 & 0.66 & 34{,}155 & 0.73 & 43.2 \\
sanders\_2020                 & 79 & 0.259 & 0.75 & 33{,}411 & 0.76 & 42.9 \\
piketty\_book                 & 78 & 0.413 & 0.66 & 27{,}550 & 0.72 & 42.9 \\
norway\_actual                & 64 & 0.091 & 0.75 & 40{,}979 & 0.88 & 40.7 \\
spain\_solidarity             & 77 & 0.365 & 0.65 & 28{,}331 & 0.71 & 41.5 \\
reform\_uk\_flavour           & 45 & 0.019 & 0.79 & 28{,}952 & 0.89 & 38.1 \\
\bottomrule
\end{tabular}
\end{table}

Values reflect the State-excluded percentile statistics
(\S\ref{sec:phi}) and the symmetric $\Phi$ definition. For
\texttt{zucman\_2pct\_spend}, top-10 share is 0.54 and Gini 0.58.
The household-sector inequality reductions are strong: the median
spend-policy halves household Gini relative to baseline Tory.

\begin{figure}[t]
\centering
\includegraphics[width=0.95\linewidth]{policies_compare.png}
\caption{Policy comparison trajectories at $N=100$. Six metrics, mean
across 6 seeds. Spending-package policies cluster in the upper half
of \emph{alive} from tick~50; Reform-flavour visibly collapses.}
\label{fig:policies-compare}
\end{figure}

Three observations.

First, \texttt{zucman\_2pct\_spend} (flat 2\% above \pounds10M with a
spending package) and \texttt{zucman\_progressive\_2\_4\_8} (the same
spending package but with the stacked 2/4/8\% schedule across
\pounds10M / \pounds100M / \pounds1B) are not separable at $N=100$:
median alive ($80$ vs $80$), top-10 share ($0.54$ vs $0.54$), and
metabolic rate ($43.1$ vs $43.1$) all agree to two decimal places.
$\Phi$ differs by 0.008 (0.412 vs 0.420), within seed noise. At
default parameters the progressive structure adds nothing detectable
above the flat 2\% bracket. We investigate the mechanism in
\S\ref{sec:top-tier-no-bind} and revisit at calibrated parameters
in \S\ref{sec:calibrated-policies}.

Second, \texttt{zucman\_2pct\_hoard}, the same wealth tax with no
spending package, gives 51 alive vs \texttt{uk\_tory\_2010\_24}'s
57. The Zucman proposal evaluated as a revenue-only instrument
costs roughly six lives per hundred at this scale relative to no
policy, while raising $\Phi$ substantially (0.418 vs 0.050) by
shrinking rentier stocks. The pattern is consistent: taxing without
spending does extract from the rentier class but does not buy a life
the way the spending package does.

Third, the \texttt{reform\_uk\_flavour} prospectus (lower taxes, no
state spending in the productive sector, higher rent extraction)
produces the lowest alive count (45) and lowest closure (0.019).
This is a near-future political scenario in the UK context and
serves as the model's most extreme present-day policy projection.

\subsection{Threshold dominates rate}
\label{sec:threshold-rate}

To isolate the effect of wealth-tax design from spending, we run two
sweeps holding the spending package fixed at the
\texttt{zucman\_2pct\_spend} configuration:

\begin{description}
\item[Sweep A] (Figure~\ref{fig:sweepA}). A single-tier wealth tax with
threshold $\theta \in \{1, 3, 10, 30, 100, 300, 1000, 3000, 10000\}$
units and rate $\tau \in [0, 0.08]$. $9\times 9$ grid, 3 seeds.
\item[Sweep B] (Figure~\ref{fig:sweepB}). A three-tier wealth tax with
fixed thresholds at $\pounds$10M, $\pounds$100M, $\pounds$1B,
varying the bottom-tier rate $\tau_b \in [0, 0.05]$ and the top-tier
marginal addition $\tau_t \in [0, 0.10]$. $9\times 9$ grid, 3 seeds.
\end{description}

\begin{figure}[t]
\centering
\includegraphics[width=0.95\linewidth]{sweep_A_heatmaps.png}
\caption{Wealth-tax Sweep~A at $N=100$. Threshold (horizontal axis,
log-spaced) versus rate (vertical). Autocatalytic closure
(top-right) plateaus from threshold $\pounds$10M down to
$\pounds$100K; above $\pounds$100M the policy loses traction.}
\label{fig:sweepA}
\end{figure}

\begin{figure}[t]
\centering
\includegraphics[width=0.95\linewidth]{sweep_B_heatmaps.png}
\caption{Wealth-tax Sweep~B at $N=100$. Bottom-tier rate (horizontal)
versus top-tier rate (vertical). Both the alive and closure surfaces
are column-identical: the top-tier rate has no effect on outcomes
because, once the bottom tier is active, no agent reaches the
$\pounds$1B threshold.}
\label{fig:sweepB}
\end{figure}

Sweep~A shows a clean monotone effect of \emph{threshold}: closure at
4\% rate falls from $\Phi = 0.59$ at threshold $\pounds$10M to
$\Phi = 0.10$ at threshold $\pounds$1B. The Zucman threshold of
$\pounds$10M is approximately the highest threshold that still
delivers a viable steady state.

\subsection{The top tier does not bind: an analytical fixed point}
\label{sec:top-tier-no-bind}

Sweep~B's ``column-identical alive/closure surfaces'' is a
tautology conditional on the post-tax wealth ceiling: any tier whose
threshold exceeds the maximum agent wealth contributes
exactly zero by construction. The economically interesting question
is \emph{where} the post-tax wealth ceiling sits, as a function of
the lower-bracket rate $\tau$ and the capital-return rate $r$.
\citet{benhabib2011distribution} provide the rigorous analytical
treatment of this question in heterogeneous-agent macro with
bequests; we recover their qualitative structure below within
the ABM. The empirical UK wealth-share targets in
\S\ref{sec:pareto} are from the WID database, in the tradition of
\citet{saez2016wealth}.

For a single rentier with wealth $w \geq \theta$, the dynamics under
capital return $r$ and a flat marginal wealth tax $\tau$ above
threshold $\theta$ are
\begin{equation}
w_{t+1} \;=\; w_t + r w_t \;-\; \tau (w_t - \theta).
\end{equation}
The fixed point $w^*$ satisfies $r w^* = \tau (w^* - \theta)$, i.e.\
\begin{equation}
w^* \;=\; \frac{\tau\, \theta}{\tau - r}
\qquad \text{when } \tau > r,
\end{equation}
and the dynamics diverge when $\tau \leq r$. Wealth tax must exceed
the rate of capital return to produce a finite stationary cap. Below
that point, agent wealth grows without analytical bound; the
empirical cap is set elsewhere, by mortality, lifespan, and the
deposit pool's drain on rentier surplus.

Figure~\ref{fig:fixed-point} compares this prediction to simulated
top wealth across the $(\tau, \theta)$ plane at $N=1000$. The
analytical prediction $w^* = \tau \theta / (\tau - r)$ tracks the
simulated top wealth quantitatively for $\tau > r$ at $\theta =
\pounds 10M$ (predicted 300, observed 308 at $\tau=0.06$;
predicted 200, observed 246 at $\tau=0.08$); at $\theta = \pounds 100M$
the simulated tops sit \emph{below} the analytical prediction
because secondary sinks (deposit pool, subsistence, inheritance tax)
act as an additional effective tax that the simple two-term dynamic
ignores. For $\tau < r$, the simulated tops are bounded only by
mortality, decreasing in $\tau$ but never reaching a finite analytical
fixed point.

\begin{figure}[t]
\centering
\includegraphics[width=0.95\linewidth]{fixed_point.png}
\caption{Analytical fixed point $w^* = \tau \theta / (\tau - r)$ vs
observed top wealth across $(\tau, \theta)$, $N=1000$, $r=0.04$,
4 seeds. The dashed line is the analytical prediction (defined only
for $\tau > r$); the solid line is the observed top-wealth median
with IQR. Secondary sinks (deposits, subsistence, inheritance)
pull the observed top below the simple analytical prediction at
high threshold.}
\label{fig:fixed-point}
\end{figure}

\paragraph{Reading.} For the wealth-tax policies modelled,
Zucman's 2\% above \pounds10M has $\tau < r = 0.04$ in our default
calibration. The simple analytical model therefore predicts
unbounded top wealth, and the simulation's observed cap
($\approx \pounds 270$M under Zucman) is set by mortality and
secondary sinks rather than by the tax-vs-compounding fixed point.
The higher Zucman tier (4\% above \pounds 100M) crosses the
$\tau > r$ threshold (the tier is inactive when $\tau = r$).
Were the simulated population ever to reach the \pounds 100M
tier, the stacked $\tau = 0.06$ above \pounds 100M would cap
them at $w^* = 0.06 \cdot \pounds 100\text{M} / (0.06 - 0.04)
= \pounds 300\text{M}$ at $r = 0.04$. At default parameters
they do not reach it, because the \pounds 10M bracket plus
secondary sinks has already held them at \pounds 270M. The
higher bracket is idle not by
tautology but because the lower bracket plus background dynamics
have already enforced a cap below it. The empirical flat-vs-
progressive equivalence at default parameters is a constructive
consequence of $\tau$-vs-$r$ competition rather than a counting
artefact, and the equivalence breaks at calibrated parameters
where the upper tiers actively bind (\S\ref{sec:calibrated-policies}).

A practical policy implication: a wealth-tax design with all tiers
at rates $\tau < r$ relies entirely on mortality and secondary
sinks to cap top wealth, which is a fragile basis. Schemes whose
\emph{lowest} tier rate exceeds the capital-return rate are
structurally robust because they cap the top analytically rather
than by accident of background dynamics. By that test the Zucman
proposal at 2\% is too low; Warren's 2\% and 6\% are also too low at
the lower thresholds and only the highest tier crosses $r$. Sanders'
top tier at 8\% comfortably exceeds typical $r$ values but its
lowest tier at 1\% does not. None of the modelled policies pass
the $\tau > r$ test at their lowest bracket; the model recovers
the long-standing concern that nominal wealth taxes are typically
below the rate of return.

\subsection{Scaling to \texorpdfstring{$N=10{,}000$}{N=10,000}: the billionaire class}
\label{sec:scale-up}

The $N=100$ analysis is open to the objection that the billionaire
class cannot statistically express itself in a 100-agent population.
We therefore scaled the model to $N=10{,}000$ for a focused subset of
eight policies. Table~\ref{tab:policies-10k} reports the
billionaire-tier wealth counts, alive, and closure. Figure
\ref{fig:billions} visualises the counts above the $\pounds$10M,
$\pounds$100M, and $\pounds$1B thresholds.

\begin{table}[t]
\centering
\small
\caption{Wealth-tier counts at $N=10{,}000$, \textbf{10 seeds}, 180
ticks (revision: previously 3 seeds). State agent excluded from
top-share statistics (cross-cutting fix). Counts are mean number of
agents whose terminal wealth (in $\pounds$) exceeds the column
threshold.}
\label{tab:policies-10k}
\begin{tabular}{lrrrrrr}
\toprule
policy & alive & $\Phi$ & top-1 share & $\geq \pounds$10M & $\geq \pounds$100M & $\geq \pounds$1B \\
\midrule
uk\_tory\_2010\_24            & 6{,}268 & 0.118 & 0.183 & 1{,}540 & 796 & \textbf{19} \\
uk\_labour\_2024              & 6{,}235 & 0.288 & 0.167 & 1{,}435 & 453 & 1 \\
zucman\_2pct\_hoard           & 6{,}103 & 0.438 & 0.106 & 1{,}440 & 27  & 1 \\
zucman\_2pct\_spend           & 8{,}053 & 0.440 & 0.095 & 1{,}762 & 69  & 1 \\
zucman\_progressive\_2\_4\_8 & 8{,}059 & 0.446 & 0.087 & 1{,}761 & 64  & 1 \\
warren\_2020                  & 7{,}699 & 0.310 & 0.103 & 1{,}749 & 704 & 1 \\
piketty\_book                 & 7{,}777 & 0.432 & 0.135 & 1{,}456 & 108 & 1 \\
reform\_uk\_flavour           & 5{,}334 & 0.060 & 0.172 & 1{,}632 & 768 & \textbf{97} \\
\bottomrule
\end{tabular}
\end{table}

Tory generates 19 billionaires in a 10{,}000-agent economy
(roughly the UK ratio of approximately 50 billionaires per 67
million population); Reform generates 97 and kills nearly half the
population. Every wealth-tax policy at default parameters collapses
the billionaire count to one. Top-1 share at \texttt{zucman\_2pct\_spend}
is 0.10, household-sector inequality compressed substantially relative
to Tory's 0.18.

\paragraph{Equivalence claim, properly tested.} The flat versus
progressive Zucman comparison at 10 seeds is
$(\text{alive}, \Phi)$ = $(8053, 0.440)$ vs $(8059, 0.446)$,
which falls within seed noise (IQRs overlap completely). The
two schedules are not separable at this scale and at default
parameters; at calibrated parameters they separate cleanly
(\S\ref{sec:calibrated-policies}, Table~\ref{tab:policies-calibrated}).

\begin{figure}[t]
\centering
\includegraphics[width=0.95\linewidth]{billionaire_count_n10k.png}
\caption{Counts of agents above wealth thresholds at $N=10{,}000$,
mean across 10 seeds. The $\geq\pounds$1B count varies from 1 under
any wealth-tax policy to 19 under the Tory baseline to 97 under the
Reform prospectus.}
\label{fig:billions}
\end{figure}

\begin{figure}[t]
\centering
\includegraphics[width=0.95\linewidth]{top_wealth_distribution_n10k.png}
\caption{Wealth of the top 200 agents by policy, log--log scale,
$N=10{,}000$. Reform-flavour and Tory baseline both show fat tails
extending past $\pounds$1B; wealth-tax policies bring the tail
sharply downward at the $\pounds$10M--$\pounds$100M region.}
\label{fig:top200}
\end{figure}

The qualitative findings of \S\ref{sec:policy} and
\S\ref{sec:top-tier-no-bind} survive scaling. The Tory baseline
generates 19 billionaires in a 10{,}000-agent economy, a rough match
to the observed UK ratio of around 50 resident billionaires per
67 million population. The Reform prospectus generates 97
billionaires and kills 47\% of the population. Every modelled
wealth-tax policy, Zucman flat, Zucman progressive, Warren,
Piketty, collapses the billionaire count to 1.

The Zucman flat-vs-progressive equivalence is not separable at 10
seeds:
$(\text{alive}, \Phi, \geq\pounds 100\text{M count})$ for the two
schedules are $(8053, 0.457, 69)$ vs $(8059, 0.464, 64)$. The
$\Phi$ difference of 0.007 sits well within seed noise. At default
parameters the top tier does not bind, for the analytical reasons
derived in \S\ref{sec:top-tier-no-bind}. At calibrated parameters
the conclusion changes; see \S\ref{sec:calibrated-policies}.

\subsection{Calibration check: the Pareto tail}
\label{sec:pareto}

Figure~\ref{fig:pareto} plots the top-200 wealth of each policy at
$N=10{,}000$ on log--log axes, with an approximate power-law fit to
the UK Sunday Times Rich List 2024 overlaid as a dashed reference.
The UK fit is $W(\text{rank}) \approx 37 \cdot \text{rank}^{-0.706}$
billion pounds, anchored to a top wealth of roughly $\pounds$37B and
approximately 165 sterling billionaires.

\begin{figure}[t]
\centering
\includegraphics[width=0.95\linewidth]{pareto_with_uk_overlay.png}
\caption{Top-200 wealth by rank for each policy at $N=10{,}000$, 3
seeds, terminal tick 180. Dashed black: approximate UK Sunday Times
Rich List 2024 power-law fit (anchored to 165 billionaires and a top
wealth of roughly $\pounds$37B). The model under-fattens the
extreme top tail relative to UK reality but produces a much higher
per-capita billionaire density. See discussion in
\S\ref{sec:pareto}.}
\label{fig:pareto}
\end{figure}

The figure makes three calibration facts visible.

First, all policy lines are shallower than the UK fit. Linear fits of
$\log W$ on $\log \text{rank}$ for the top 50 agents give Pareto
slopes around $-0.07$ to $-0.26$, compared to the UK fit slope of
$-0.706$. The model's top tail is markedly flatter than the empirical
distribution: many similar-wealth agents at the top rather than a few
extreme outliers.

Second, the maximum wealth in the Tory baseline is approximately
$\pounds$4 to $\pounds$6B at the median seed, compared to roughly
$\pounds$37B in the UK list. The model under-produces the very-top
fortunes by approximately one order of magnitude.

Third, and weighing against the first two, the model
\emph{over-produces} the lower-billionaire region: 19 billionaires per
$10{,}000$ agents under Tory is about $2500\times$ the UK per-capita
rate of around 50 billionaires per 67 million. We have too
many billionaires per capita, but each one is too poor.

This combination, a flatter and lower top tail with too many
people in the billionaire band, points to a small number of model
features that are jointly responsible. Capital return at 4\% per tick
likely undershoots the genuine compounding of top dynastic wealth.
Rentier proportion fixed at 3\% of the population is plausibly an
order of magnitude too high for the strict billionaire class but
fine for the broader $\pounds$10M$+$ class. The bandwidth-coupled
inheritance mechanism may also redistribute too much knowledge between
generations to allow the kind of multi-generation dynastic
concentration the empirical tail reflects.

A partial calibration over capital return $r \in \{0.04, 0.06, 0.08,
0.10\}$, rent share $\in \{0.20, 0.35, 0.50\}$, and inheritance tax
$\iota \in \{0.0, 0.2, 0.4\}$ (180 cells, 5 seeds at $N=1000$),
fitting to UK \citep{saez2016wealth} WID targets of top-1\% share
0.30, top-10\% share 0.57, top-0.1\% share 0.13, and Pareto slope
$-0.706$, returns the best L2 fit at $r=0.06$, rent share 0.50,
inheritance tax 0.40, yielding observed top-1\% 0.15, top-10\% 0.65,
top-0.1\% 0.06, Pareto slope $-0.673$. The slope matches the UK
Rich List slope almost exactly; the absolute levels remain too low
by roughly a factor of two at top-1\% and top-0.1\%. The slope match
is illustrative rather than a fully validated calibration; this
configuration keeps only 26\% of the population alive at terminal
tick, well below UK reality. There is a Pareto frontier here between
matching empirical tail steepness and empirical mortality, and we
have not found a setting that hits both. The fact that the slope is
reachable indicates the model can produce realistic-tail behaviour
in some configuration; \S\ref{sec:calibrated-policies} reports how
much the policy conclusions move when the policy suite is rerun at
the calibrated configuration. Future work should free the population
proportions (the rentier share is hard-coded at 3\%, plausibly too
high for the strict billionaire class) and fit jointly.

\subsection{Policy results at calibrated parameters}
\label{sec:calibrated-policies}

The partial calibration in \S\ref{sec:pareto} matched the UK Pareto
slope at $r = 0.06$, rent share $0.50$, inheritance tax $0.40$. The
policy results reported in \S\ref{sec:policy} and \S\ref{sec:scale-up}
are at the default $r = 0.04$, rent share $0.35$. The
configuration that earns the \pounds\ axis is not the one in which
the policy comparisons were made. This section corrects that by
rerunning the policy suite with $r = 0.06$ and rent share $0.50$
held fixed across policies, leaving each policy's other dial
settings unchanged. Table~\ref{tab:policies-calibrated} reports
$N=10{,}000$, 6 seeds, 180 ticks.

\begin{table}[t]
\centering
\small
\caption{Policy results at calibrated parameters ($r=0.06$, rent
share $=0.50$), $N=10{,}000$, 6 seeds, 180 ticks. State excluded
from top-1 share. Compare to Table~\ref{tab:policies-10k}
(default parameters). Wealth thresholds in \pounds.}
\label{tab:policies-calibrated}
\begin{tabular}{lrrrrrr}
\toprule
policy & alive & top-1 share & $\geq$\pounds10M & $\geq$\pounds100M & $\geq$\pounds1B & max wealth \\
\midrule
uk\_tory\_2010\_24 & 2{,}230 & 0.198 & 866 & 309 & \textbf{300} & \pounds48B \\
uk\_labour\_2024 & 3{,}029 & 0.157 & 1{,}360 & 800 & 262 & \pounds12B \\
zucman\_2pct\_hoard & 2{,}532 & 0.262 & 910 & 306 & 10 & \pounds4.1B \\
zucman\_2pct\_spend & 7{,}734 & 0.233 & 1{,}746 & 801 & 46 & \pounds2.4B \\
zucman\_progressive\_2\_4\_8 & 7{,}744 & 0.143 & 1{,}750 & 720 & \textbf{0} & \pounds0.4B \\
warren\_2020 & 7{,}234 & 0.179 & 1{,}740 & 804 & 140 & \pounds1.5B \\
piketty\_book & 7{,}363 & 0.230 & 1{,}441 & 800 & 20 & \pounds1.7B \\
reform\_uk\_flavour & 2{,}526 & 0.181 & 1{,}010 & 341 & \textbf{324} & \pounds77B \\
\bottomrule
\end{tabular}
\end{table}

Three things move materially. First, billionaire counts grow under
every policy except the progressive Zucman: Tory now produces 300
billionaires (vs 19 at default), Reform produces 324 (vs 97), and
even Zucman 2\% flat with spending leaves 46 (vs 1). The model at
calibrated parameters is much closer to the UK Rich List
distribution (Tory max wealth \pounds48B vs the UK Rich List top of
roughly \pounds37B) but the per-capita billionaire density now
\emph{over}-shoots in the other direction.

Second, mortality is much higher in non-spend policies under
calibrated parameters: Tory keeps 22\% of the population alive at
$r=0.06$, vs 63\% at $r=0.04$. The combination of higher capital
return and higher rent share extracts faster than the population
can replace itself absent a spending package. This is the
``calibrated config kills three-quarters'' problem the previous
section flagged, now made concrete across all policies.

Third, and most importantly,
\textbf{Zucman progressive now binds the top tier and Zucman flat
does not}. Under calibrated parameters Zucman progressive produces
\emph{zero} billionaires while Zucman flat leaves 46, a clean
separation that did not exist at default $r=0.04$. Top-1 share
also separates clearly: 0.143 progressive vs 0.233 flat. The
prior draft's ``top tier does not bind'' result was specific to
the default-parameter regime where wealth never reaches the
\pounds1B bracket; once $r$ is raised to the calibration target, the
upper tier becomes the decisive instrument.

\paragraph{Implications for the analytical fixed point.} At
calibrated $r = 0.06$, none of the policy lowest-tier rates exceed
$r$ (Zucman 2\%, Warren 2\%, Piketty 1\% at lower brackets, Sanders
1\%). The cap-by-mortality regime identified in
\S\ref{sec:top-tier-no-bind} now bites harder, which we observe as
the high mortality of the non-spend cells. The Zucman progressive
top tier of 4\% above \pounds100M (stacking to 8\% above \pounds1B)
is the one rate among the modelled policies that exceeds the
calibrated capital return, and it is the one policy that eliminates
billionaires. The \pounds1B tier is no longer a tautology; it is the
mechanism by which the policy succeeds in this regime. The
$\tau$-vs-$r$ design test is preserved at calibrated parameters
and becomes more rather than less significant.

\paragraph{Implications for the policy contribution.} A practical
political reading. The Zucman headline of ``2\% on the
ultra-wealthy'' is robust as a household-inequality compressor
under either parameter regime, but it does not eliminate
billionaires under realistic capital-return dynamics; only the
progressive form does. The political case for the upper tiers
(currently presented in Zucman's proposal as a sign of seriousness
rather than as the load-bearing instrument) is, under our
calibration, the central mechanism. Conversely, if the genuine
capital return is closer to the default $r=0.04$ than to the
calibration target $r=0.06$, the lower tier alone suffices and the
progressive structure is symbolic. Either way the \pounds10M
threshold is doing real work.

\subsection{Open-ended skill space: a minimal demonstration}
\label{sec:open-skills}

The cultural model presented above operates in a fixed K=5 skill space.
No skill dimension that did not exist at $t=0$ can ever appear, which
is a category limitation relative to the open-ended-evolution
tradition \citep{taylor2019evolutionary, banzhaf2016defining}. We
address this with a minimal demonstration extension. In a stand-alone
simulation the skill vector
is variable-length: at every reproduction event, with probability
proportional to $\log(w_{\text{parent}})$, a previously-undiscovered
skill dimension enters the global pool and is initialised on the child
to a fresh Beta(2,2) value (subsequent inheritance follows the same
bandwidth rules as the main model).  Five hundred initial agents, 400
ticks, five seeds.

\begin{figure}[t]
\centering
\includegraphics[width=0.95\linewidth]{open_skills.png}
\caption{Open-ended skill space demonstration. \textbf{Left:} active
skill-pool size over time. The closed variant (grey, novelty
rate 0) stays at five; the open variant (blue, novelty
proportional to $\log(\text{parent wealth})$) climbs from five to
about ten.  \textbf{Middle:} per-agent breadth grows accordingly.
\textbf{Right:} the discoverers of new skills: most novelty events
fire on worker parents (because workers are the most numerous class),
but the upper tail of the discoverer-wealth distribution sits at
$10^{4}$ in model units, indicating that high-wealth parents do
contribute to skill discovery disproportionately to their numbers.}
\label{fig:open-skills}
\end{figure}

The extension does not feed back into the policy results of the main
paper, and we do not claim it as a load-bearing finding. It exists to
show that the model's core mechanism can be extended in a TAP-style
adjacent-possible direction \citep{hordijk2023emergence} without
fundamental redesign, and to demonstrate that the wealth-coupled
discovery channel concentrates novelty production in wealthier
lineages on a per-capita basis. The per-capita discovery rate (right
panel) is 0.012 per agent-generation for RENTIER, 0.009 for WORKER,
and 0.005 for MAKER; rentiers discover at $1.3\times$ the worker rate
and $2.4\times$ the maker rate, despite being far less numerous in
absolute terms. The wealth-coupled discovery channel does concentrate
novelty in wealthier lineages on a per-capita basis as well as in
absolute count.

A more substantive treatment, with open-ended novelty integrated
into the full economic model and the policy results re-run under
endogenous skill discovery, is the natural follow-up paper; we cite
\citet{taylor2019evolutionary} and \citet{banzhaf2016defining} as
the standard references for that direction.

\subsection{The regressive trap was state contamination}
\label{sec:regressive-trap}

An earlier version of this analysis reported that at $N=10{,}000$
with a fixed spending package, top-1\% share \emph{increases} from
$0.37$ to $0.84$ as the wealth-tax rate rises from zero to 8\% at
threshold \pounds 1M. That finding included the State agent in the
wealth distribution, so as the wealth tax shifted revenue from
private households to the State, the State's growing share dominated
the percentile statistics. With the State excluded and a
mortality-corrected variant computed alongside, the picture changes
materially. We report top-1\% share over the rate axis for the seven
threshold levels (thresholds in model units, where 1 unit =
\pounds 100k):

\begin{table}[t]
\centering
\small
\caption{Top-1\% share at $N=10{,}000$, 2 seeds, 150 ticks, State
excluded. Each row is a fixed threshold; each column a rate from 0
to 0.08. The mortality-corrected and raw values agree to within
0.005 throughout; we report raw values here.}
\label{tab:taxsweep10k}
\begin{tabular}{rrrrrrrr}
\toprule
$\theta$ (\pounds) & $\tau=0$ & 0.013 & 0.027 & 0.040 & 0.053 & 0.067 & 0.080 \\
\midrule
1M    & 0.18 & 0.14 & 0.11 & 0.09 & 0.07 & 0.06 & 0.06 \\
3M    & 0.18 & 0.13 & 0.09 & 0.07 & 0.05 & 0.04 & 0.04 \\
10M   & 0.18 & 0.12 & 0.08 & 0.06 & 0.05 & 0.04 & \textbf{0.04} \\
30M   & 0.18 & 0.11 & 0.08 & 0.07 & 0.06 & 0.05 & 0.05 \\
100M  & 0.18 & 0.12 & 0.10 & 0.09 & 0.08 & 0.08 & 0.08 \\
300M  & 0.18 & 0.16 & 0.14 & 0.14 & 0.13 & 0.13 & 0.12 \\
1B    & 0.18 & 0.18 & 0.18 & 0.18 & 0.18 & 0.17 & 0.17 \\
\bottomrule
\end{tabular}
\end{table}

\begin{figure}[t]
\centering
\includegraphics[width=0.95\linewidth]{tax_sweep_n10k.png}
\caption{Wealth-tax sweep at $N=10{,}000$, State excluded. Threshold
on horizontal log-spaced axis (\pounds 1M to \pounds 1B); rate on
vertical. The top-1\% share surface is monotone decreasing in rate
at every threshold; the apparent ``regressive trap'' of an earlier
analysis was a state-contamination artefact and disappears under
the corrected metric.}
\label{fig:taxsweepN10k}
\end{figure}

There is no regressive trap in this model once the State is
excluded from the wealth distribution. Top-1\% share decreases
monotonically with the tax rate at every threshold; the strongest
compression is at the \pounds10M Zucman threshold (top-1\% falls
from 0.18 to 0.04, a $4.4\times$ reduction). The mortality
correction itself moves the metric by less than 0.005 across the
entire grid, indicating that excluding the State (rather than
correcting for excess mortality) is what changes the answer.

\paragraph{The substantive finding that does survive.} The threshold
ranking holds: top-1\% reduction is most aggressive at intermediate
thresholds (\pounds 10M to \pounds 30M) and tapers off at higher
thresholds. The Zucman \pounds10M threshold turns out to compress
the top quintile most efficiently of any threshold we tested. At
\pounds 1B threshold, the tax effectively does nothing (the
threshold catches almost no agents). This is the genuine
\emph{progressive} finding the original draft mis-stated as
``regressive''; corrected, it points the same direction the
political case for Zucman points.

%==========================================================================
\section{Discussion}
\label{sec:discussion}
%==========================================================================

\subsection{What we believe is novel}

The class-specific skill-to-wealth elasticity decomposition
(\S\ref{sec:culture}, Table~\ref{tab:elasticity}) is the freshest
contribution. The cultural-economic asymmetry the bandwidth
coupling generates is not a structural artefact of how the model is
wired; it is a consequence of one half of the economy (extractive
classes) having near-zero skill-to-wealth feedback while the other
half (productive classes) has strong feedback, which a regression on
40-tick wealth gain shows directly.

The analytical $w^* = \tau \theta / (\tau - r)$ fixed-point
characterisation (\S\ref{sec:top-tier-no-bind}) re-frames the
top-tier-does-not-bind result from a tautology conditional on the
post-tax ceiling
into a derivable consequence of tax-vs-compounding competition.
\citet{benhabib2011distribution} provide the rigorous heterogeneous-
agent macro treatment of bequest dynamics and wealth-distribution
tails; we recover the qualitative $\tau$-vs-$r$ structure in the
ABM setting.

The (tax $\times$ spending $\times$ inheritance) clean factorial
(\S\ref{sec:state-spending}, Table~\ref{tab:factorial}) disaggregates
three policy levers. Spending raises alive counts without
moving $\Phi$; tax-rate schedules compress household-sector
inequality; inheritance taxation is the most $\Phi$-effective
single lever.

The wealth-coupled cultural-transmission bandwidth mechanism
\citep{henrich2001evolution, henrich2004demography,
spangler2022distribution} is exercised across the relevant wealth
range under the log-wealth coupling (\S\ref{sec:culture}).

The $\Phi$ flow-share metric (\S\ref{sec:phi}) is defined
symmetrically across productive and extractive accrual. We do not
claim it is an autocatalytic-set property in the
\citet{hordijk2004detecting} RAF sense; it is a per-tick flow ratio
operationalising the productive-vs-extractive distinction without
implementing the topological RAF property.

\subsection{Honest accounting of where the model still falls short}

The household-sector inequality reductions under spend policies are
robust across all metrics we report, but several quantitative claims
remain conditional on parameter choices. The ``top tier does not
bind'' result holds at default parameters; at calibrated parameters
(\S\ref{sec:calibrated-policies}) the top tier binds clearly and the
progressive Zucman schedule outperforms the flat. Which result
should be considered the headline depends on whether the calibrated
configuration is judged a fair representation of UK conditions; we
present both and discuss the trade-off rather than picking one.

The original analysis reported a ``regressive trap'' (top-1\% share
rising with the tax rate at low thresholds) that was the result of
including the State agent in the wealth distribution; it does not
survive State exclusion. We retract that finding.

The Pareto-tail calibration is partial. The slope is reachable at
$r=0.06$, rent share 0.50, inheritance tax 0.40, but at the cost of
mortality that the model cannot reproduce on UK demographic
benchmarks. We do not yet have a fitted scenario that matches both
slope and level and survival simultaneously.

\subsection{What the model cannot do}

The most consequential remaining limitations: no spatial structure;
no firm formation; no fiscal-crisis dynamics; toy speculator
dynamics that do not produce bubbles; a closed K=5 skill space in
the main results (a minimal open-skills demonstration in
\S\ref{sec:open-skills} shows how the mechanism extends to
open-ended skill discovery but does not fold back into the policy
analysis); a single lender pool with no shadow-banking.
\citet{charbonneau2021fidelity}'s ``grain problem'' for fidelity
measures applies to our K=5 skill discretisation and we acknowledge
it as a methodological choice rather than a natural quantity.

\subsection{Positioning for publication}

\begin{description}
\item[JASSS] (Journal of Artificial Societies and Social Simulation).
Lead with (i) the productive-flow-share metric $\Phi$ and the
Sankey visualisation (Figure~\ref{fig:sankey}), which is the
single best artefact for explaining the metric to a non-technical
audience; (ii) the phase portrait under state-excluded Gini
(Figure~\ref{fig:phase}, Table~\ref{tab:regimes}); and (iii) the
clean (tax $\times$ spending $\times$ inheritance) factorial
(\S\ref{sec:state-spending}, Table~\ref{tab:factorial}), which
shows that alive, Gini, and $\Phi$ are three distinct policy
targets answering to three distinct levers. The political framing
reads natively in JASSS. For an ALIFE submission a more conservative
title, e.g.\ replacing ``autocatalytic-closure'' with
``productive-flow share,'' would reduce friction with the
RAF-literate audience while preserving the conceptual lineage.

\item[JEIC] (Journal of Economic Interaction and Coordination).
Lead with the analytical $\tau$-vs-$r$ fixed point
(\S\ref{sec:top-tier-no-bind}), the wealth-tax design factorisation
under both default and calibrated parameters
(\S\ref{sec:calibrated-policies}), and the calibration-vs-policy
sensitivity result. Direct engagement with
\citet{spangler2022distribution} and \citet{cardoso2025wealth}
positions the paper; \citet{benhabib2011distribution} is the
analytical neighbour to engage on the fixed-point side.

\item[Artificial Life / ALIFE]. Lead with the class-specific
skill-to-wealth elasticity decomposition (Table~\ref{tab:elasticity}),
the open-skill-space demonstration (\S\ref{sec:open-skills}) with
per-capita discovery rate, and the bandwidth decomposition
(Figure~\ref{fig:bandwidth}). The framing is cultural erosion under
economic pressure rather than cumulative cultural evolution. We
would retitle around ``productive-flow share'' here.
\end{description}

\subsection{Honest reading}

The four findings the model produces that we did not anticipate are:
the analytical $\tau$-vs-$r$ fixed point that determines the
post-tax wealth cap (\S\ref{sec:top-tier-no-bind}); the class-specific
skill-to-wealth elasticity decomposition that grounds the
cultural-economic asymmetry in mechanism rather than architecture
(\S\ref{sec:culture}); the dis-conflation of the spending,
tax-schedule, and inheritance levers in the clean factorial; and the
result that at calibrated capital return $r=0.06$ the Zucman
progressive schedule sharply outperforms the flat schedule (zero
billionaires vs 46), reversing the default-parameter equivalence
(\S\ref{sec:calibrated-policies}).

The two known-in-spirit findings are: state spending keeps people
alive (a kinetic-ABM result, attributable to
\citet{bertotti2018modanese} and \citet{cardoso2025wealth}); and
inheritance reform requires redistribution
(\citealp{piketty2013theory} in analytical form).

If the work has a thesis it is: \emph{the productive-flow share
$\Phi$, the household-sector Gini, and the population-alive count
are three distinct policy targets that respond to distinct policy
levers; any political framing of a wealth-tax design as ``the
wealth tax'' lumps together effects that should be evaluated
separately.} Inheritance taxation raises $\Phi$; tax-rate schedules
compress household inequality; spending keeps people alive.
Quantitative policy claims expressed in \pounds are partially
supported by the calibration in \S\ref{sec:pareto} but rest on a
configuration that does not simultaneously match the empirical
mortality; threshold-comparison results are also reported in
dimensionless terms (fraction of the wealth-distribution
median) wherever the pound label was load-bearing, with the partial
calibration result in \S\ref{sec:pareto} backing the
\pounds10M anchor in particular.

%==========================================================================
\section{Conclusion}
%==========================================================================

We have presented an agent-based model of wealth inequality whose
distinguishing contributions are a productive-flow share metric
for flow-based system viability and a wealth-coupled bandwidth
mechanism for skill transmission. Across thirteen stylised
scenarios, eleven real-world wealth-tax policy specifications, a
clean three-factor factorial, and a partial calibration to the UK
Sunday Times Rich List Pareto slope, we have identified the eight
findings enumerated in the abstract: three robust regimes;
three-distinct-levers factorisation of policy effects; class-specific
skill-to-wealth elasticities; inheritance trapdoor; non-separability
of flat-vs-progressive Zucman at default parameters with clean
separation at calibrated parameters; full retraction of the
``regressive trap'' as state contamination; billionaire counts at
$N=10{,}000$; and the calibration partial fit.

The model's most consequential political claim, post-revision, is
that wealth-tax design is not a single choice but a triple of
choices, each addressing a different metric. State spending raises
alive counts; tax-rate schedules compress household-sector inequality;
inheritance taxation raises the productive-flow share. Within the
wealth-tax-rate choice itself, the relationship between the
marginal rate $\tau$ above the lowest threshold and the rate of
capital return $r$ determines whether the policy analytically caps
top wealth (when $\tau > r$) or leaves the cap to be set by
secondary dynamics (mortality, lifespan, deposit drain). None of
the eleven real-world proposals modelled have $\tau > r$ at their
lowest tier under default $r=0.04$, which puts them all in the
``cap-by-mortality'' regime that the regressive-trap section
quantifies. A wealth-tax design with its lowest rate above the
capital-return rate is structurally robust; designs with lower rates
rely on mortality to cap accumulation, which is a fragile basis.

Code and data are available at the project repository; reproduction
requires Python 3.12 with NumPy and Matplotlib.

\bibliographystyle{plainnat}
\bibliography{references}

%==========================================================================
\appendix
%==========================================================================

\section{ODD-style summary}
\label{app:odd}

A brief ODD (Overview, Design concepts, Details) summary follows;
a full ODD+D specification is deferred to a separate online supplement.

\paragraph{Overview.}
\textbf{Purpose.} Make qualitative wealth-inequality dynamics
visible under a panel of policy levers, including a metric for the
productive-vs-extractive accrual balance.
\textbf{Entities.} Seven agent classes
(WORKER, MAKER, EMPLOYER, LENDER, RENTIER, SPECULATOR, STATE) at
fixed proportions per 100 agents (72/9/7/5/3/3/1).
\textbf{State variables.} Each agent has wealth, debt, age,
lifespan, a five-dimensional skill vector, alive flag, peak
lifetime wealth, bandwidth dimensions at birth.
\textbf{Scales.} One tick is one stylised year. The base run is
$N=100$ agents at 250 ticks. Scale-up to $N=10{,}000$ uses
population\_scale = 100 (state stays singleton; per-capita stays
invariant). One model wealth unit corresponds to \pounds100k.

\paragraph{Process overview and scheduling.} Per tick, in fixed
order: (1) Production --- makers and employers hire workers from a
skill-sorted queue and pay wages; makers produce output at multiplier
$M$; employers capture markup; the state optionally hires
unemployed workers and purchases maker output;
(2) Extraction --- workers pay rent to rentiers in proportion to
rentier wealth; rentiers earn capital return $r$; speculators
update by Brownian increment with drift and volatility;
(3) Consumption and credit --- every non-state agent pays
subsistence; shortfalls are borrowed from the lender pool, with
allocation proportional to lender capacity; death by starvation
if no capacity remains;
(4) Interest and repayment --- debtors pay interest and (when
liquid) repay principal; debt written off above ceiling;
(5) Lender deposits --- a fraction of surplus wealth from rentier,
maker, employer agents is deposited with lenders; optional state
recapitalisation to a configured floor;
(6) Tax and UBI --- progressive flat-plus-tiered wealth tax to
state; UBI dispersed to non-state agents from state wealth;
(7) Life --- agents whose age exceeds lifespan die naturally; a
child of the same class is born with inherited wealth (after
inheritance tax and cap) and bandwidth-limited inherited skills;
(8) Bookkeeping --- update peak-wealth, recent-deaths log,
metrics; prune dead agents every ten ticks.

\paragraph{Design concepts.}
\emph{Emergence.} Wealth distribution, regime classification, and
cultural-breadth distribution emerge from the dial settings.
\emph{Adaptation.} Agents do not adapt within their lifetime;
class-typical behaviour is hard-coded; behavioural variability
comes from class-typical skill biases and stochastic lifespans.
\emph{Sensing.} Limited; workers do not choose employers; agents
do not perceive prices or expectations.
\emph{Interaction.} Mediated through queues (hiring),
proportional shares (lender pool, rentier pool), and the state.
\emph{Stochasticity.} Per-agent lifespan, debt at birth (for
workers), speculator returns, skill mutation, and seed-controlled
initialisation are stochastic.
\emph{Collectives.} Implicit only --- agents do not form groups
beyond their class label.
\emph{Observation.} The metrics reported per tick: Gini (default
state-excluded), $\Phi$, alive, total wealth, total debt,
cultural breadth and diversity, knowledge lost, births and deaths
counters, top-1 and top-10 shares (raw and mortality-corrected),
state wealth, wealth by type.

\paragraph{Details.}
\emph{Initialisation.} Class shares lognormally distributed within
class; ages staggered uniformly across lifespans; workers given
random small debt with probability 0.18.
\emph{Input data.} None. All parameters are dials.
\emph{Submodels.} See \S\ref{sec:model} for production, extraction,
credit, taxation, transmission. Dials documented in
Table~\ref{tab:dials}.

\section{Dial reference}
\label{app:dials}

Default values for the dial panel are listed in Table~\ref{tab:dials}.

\begin{table}[h]
\centering
\small
\caption{Dial panel and default values. Tax thresholds are in model
units; one unit corresponds to $\pounds$100{,}000.}
\label{tab:dials}
\begin{tabular}{lrl}
\toprule
dial & default & group \\
\midrule
wage                          & 1.0 & production \\
maker\_multiplier             & 1.6 & production \\
employer\_markup              & 0.4 & production \\
workers\_per\_maker           & 2   & production \\
workers\_per\_employer        & 6   & production \\
capital\_return ($r$)         & 0.04 & capital \\
rent\_share\_of\_wage         & 0.35 & capital \\
interest\_rate ($i$)          & 0.06 & capital \\
spec\_drift / spec\_vol       & 0.015 / 0.12 & capital \\
flat\_tax\_rate               & 0.0 & state \\
tax\_progressivity            & 0.0 & state \\
wealth\_tax\_threshold        & 50 ($=\pounds$5M) & state \\
wealth\_tax\_tiers            & ()  & state \\
ubi                           & 0.0 & state \\
state\_employs\_fraction      & 0.0 & state \\
state\_employs\_wage          & 1.0 & state \\
state\_buys\_maker\_output    & 0.0 & state \\
state\_recap\_lenders         & 0.0 & state \\
subsistence ($\sigma$)        & 0.6 & life \\
debt\_ceiling                 & 25  & life \\
lifespan\_mean / std          & 70 / 12 & life \\
starvation\_lifespan\_penalty & 20  & life \\
inheritance\_cap              & $\infty$ & transmission \\
inheritance\_tax\_rate        & 0.0 & transmission \\
memory\_bandwidth\_base ($b_0$) & 0.5 & transmission \\
memory\_bandwidth\_wealth\_scale ($b_w$) & 0.5 & transmission \\
skill\_mutation\_rate         & 0.05 & transmission \\
population\_scale             & 1 & logistics \\
seed                          & 42 & logistics \\
\bottomrule
\end{tabular}
\end{table}

\section{Scenario specifications}
\label{app:scenarios}

The thirteen scenarios at $N=100$ are obtained from
\texttt{current\_world} (all dials at defaults) by overriding the
non-default dials specified in \texttt{scenarios.py} in the project
repository. Distinctive non-default settings include:
\texttt{stevenson\_world} (rent\_share = 0.55, $r$ = 0.05, $i$ = 0.09,
markup = 0.55); \texttt{maker\_economy} ($M$ = 2.2, rent = 0.15,
$r$ = 0.01, markup = 0.2); \texttt{nordic} (UBI = 0.6, state employs
0.6, state buys 10, inheritance tax 0.4, inheritance cap 500, rent
0.25); \texttt{socialist} (UBI 0.8, state employs 0.9, state buys 15,
inheritance tax 0.8, rent 0.10); \texttt{feudal} (rent 0.6, $r$ 0.06,
markup 0.5, no inheritance tax, $b_0$ 0.2, $b_w$ 1.0).

\section{Wealth-tax policy specifications}
\label{app:policies}

The eleven policies are specified in \texttt{policies.py}. Tax tiers
are stacked-marginal form $(\theta, \tau)$ tuples, with threshold in
model units (one unit = $\pounds$100{,}000).

\begin{itemize}
\item \texttt{uk\_tory\_2010\_24}: no wealth tax; inheritance tax 0.4.
\item \texttt{uk\_labour\_2024}: flat tax 0.002, progressivity 0.005;
inheritance tax 0.5; state employs 0.1.
\item \texttt{zucman\_2pct\_hoard}: tiers (100, 0.02); inheritance
tax 0.4; no spending package.
\item \texttt{zucman\_2pct\_spend}: tiers (100, 0.02); inheritance
tax 0.4; UBI 0.4, state employs 0.3, state buys 5.0.
\item \texttt{zucman\_progressive\_2\_4\_8}: tiers (100, 0.02),
(1000, 0.02), (10000, 0.04); spending package as above.
\item \texttt{warren\_2020}: tiers (500, 0.02), (10000, 0.04);
spending package.
\item \texttt{sanders\_2020}: tiers (320, 0.01), (3200, 0.02),
(32000, 0.05); spending package.
\item \texttt{piketty\_book}: tiers (10, 0.001), (50, 0.009),
(1000, 0.01); inheritance 0.6; spending package.
\item \texttt{norway\_actual}: tiers (1.3, 0.01); no inheritance tax;
state employs 0.4, state buys 4, lender backstop 200.
\item \texttt{spain\_solidarity}: tiers (7, 0.002), (30, 0.005),
(100, 0.01), (1000, 0.015); inheritance 0.3; spending package.
\item \texttt{reform\_uk\_flavour}: no wealth tax; inheritance 0.2;
no spending package; rent 0.4.
\end{itemize}

\section{Sweep specifications}
\label{app:sweeps}

The five $N=100$ scenario sweeps were performed on $9\times 9$ grids
with three seeds each over the ranges:

\begin{itemize}
\item \texttt{rent\_vs\_r}: rent $\in [0.05, 0.65]$, $r \in [0.005, 0.06]$.
\item \texttt{tax\_vs\_ubi}: flat tax $\in [0, 0.025]$, UBI $\in [0, 1.5]$.
\item \texttt{memory}: $b_0 \in [0.2, 1.0]$, $b_w \in [0, 1]$.
\item \texttt{state\_action}: state\_employs\_fraction $\in [0, 1]$,
state\_buys\_maker\_output $\in [0, 20]$.
\item \texttt{inheritance\_x\_ubi}: $\iota \in [0, 1]$, UBI $\in [0, 1.5]$.
\end{itemize}

The two wealth-tax sweeps at $N=100$ used thresholds
$\{1, 3, 10, 30, 100, 300, 1000, 3000, 10000\}$ (single-tier) and
fixed thresholds $(100, 1000, 10000)$ (three-tier), with rates over
$[0, 0.08]$ and $[0, 0.10]$ respectively.

The scale-up wealth-tax sweep at $N=10{,}000$ used a $7\times 7$ grid
of thresholds $\{10, 30, 100, 300, 1000, 3000, 10000\}$ and rates
$[0, 0.08]$ with two seeds each, 150 ticks.

\end{document}
